% ============================================================================
% P2 Mid replacement text for the Pakalorie FYP report (Overleaf).
% Every number below was measured on the built system; see docs/P2_MID.md.
% Each block says which existing section it replaces. Plain LaTeX only
% (tabular + \hline), so it should compile in your existing class.
% Adjust \section / \subsection levels to match your template.
% ============================================================================


% ----------------------------------------------------------------------------
% REPLACES 4.4.7 "Module 7: Model Optimisation and Quantisation"
% ----------------------------------------------------------------------------
\subsubsection{Module 7: Model Optimisation and Quantisation}
\textbf{Purpose:} make the recognition model smaller so it can be downloaded
and, later, run on a phone. \textbf{Inputs:} the trained YOLOv8n classification
model in ONNX format (217 classes, 224$\times$224 input).
\textbf{Processing:} the model is quantised to 8-bit integers in two ways:
dynamic quantisation, which converts only the weights, and static quantisation,
which also converts the activations using 200 calibration images. Each version
is compared with the original for file size, CPU latency and accuracy on the
same held-out images. \textbf{Outputs:} a comparison table (Table~\ref{tab:quant})
and the recommended model file.


% ----------------------------------------------------------------------------
% REPLACES 4.4.8 "Module 8: Data Visualisation and Calorie Compensation"
% ----------------------------------------------------------------------------
\subsubsection{Module 8: Data Visualisation and Calorie Compensation}
\textbf{Purpose:} help the user see trends and keep to the daily goal after a
heavy or light day. \textbf{Inputs:} saved meals, activities logged by the user,
and the daily calorie goal. \textbf{Processing:} meals are grouped by day. The
module looks at the last three days on which the user logged meals. The surplus
(calories eaten minus the goal plus calories burned that day) is added up and
divided by three, so each day's surplus is repaid over three days. The change to
today's goal is limited to 15\% of the base goal and is never allowed to take the
goal below a safe minimum of 1200~kcal (or below the user's own goal, if that is
already lower). Calories burned by logged activity are added to today's goal.
\textbf{Outputs:} an adjusted daily goal with a one-line explanation, a disclaimer
that the numbers are estimates and not medical advice, and charts of calories,
macronutrient share, weight and water intake.


% ----------------------------------------------------------------------------
% REPLACES 4.4.9 "Module 9: Real-Time Inference Pipeline"
% ----------------------------------------------------------------------------
\subsubsection{Module 9: Real-Time Inference Pipeline}
\textbf{Purpose:} give the user a result while the camera is pointed at the food,
without pressing the shutter. \textbf{Processing:} in live mode the app captures a
small, low-quality camera frame, sends it to the server's recognition service
with the YOLOv8 engine, and shows the best guess with its confidence. Only one
frame is in flight at a time; the next frame is captured as soon as the reply
arrives, so the server latency sets the frame rate. Guesses below 35\% confidence
are labelled ``low confidence''. The live mode uses only the YOLOv8 model, so it
does not call any paid language-model service. \textbf{Outputs:} a live label,
the confidence, the round-trip time and the frames per second.
In this version the pipeline is streamed from the phone to the server; running
the quantised model on the phone itself is listed as future work.


% ----------------------------------------------------------------------------
% REPLACES 5.4.4 "Optimisation and Quantisation (Module 7)"
% ----------------------------------------------------------------------------
\subsubsection{Optimisation and Quantisation (Module 7)}
The trained model was exported to ONNX and quantised to 8-bit integers with the
ONNX Runtime quantisation tools, using a script kept in the repository
(\texttt{ml/scripts/quantize.py}). Two variants were built: dynamic quantisation
and static quantisation in QDQ format, calibrated on 200 images that were left out
of the test set. All three models were then run on 1516 held-out images on a CPU.
An FP16 version was not built. Table~\ref{tab:quant} shows the result.

\begin{table}[h]
\centering
\caption{Quantisation results (CPU, 1516 held-out images)}
\label{tab:quant}
\begin{tabular}{|l|c|c|c|c|}
\hline
\textbf{Model} & \textbf{Size (MB)} & \textbf{Latency (ms)} & \textbf{Top-1} & \textbf{Top-5} \\
\hline
Original (FP32)      & 6.89 & 2.4 & 0.700 & 0.927 \\
INT8 dynamic         & 1.79 & 6.6 & 0.592 & 0.885 \\
INT8 static (QDQ)    & 1.89 & 2.1 & 0.663 & 0.925 \\
\hline
\end{tabular}
\end{table}

Static INT8 quantisation is the version to use. It is 3.6 times smaller than the
original, as fast on this CPU, and loses 3.7 points of top-1 accuracy and 0.2 points
of top-5 accuracy. Dynamic quantisation gives the smallest file but loses 10.8
points of top-1 accuracy and is slower, because the CPU runtime has no fast
integer convolution, so it is not recommended.

\textit{Limitation.} The test images were rebuilt on a new computer from the two
Kaggle datasets with the same split settings. The trained model has 217 classes,
while the earlier training notes list 218, so some of these images may have been
seen during training. This is why the original model scores 0.700 here but 0.585
top-1 (0.866 top-5) in the official evaluation of Section~6.4. Use Table~\ref{tab:quant}
to compare the three versions with each other, and cite the official figures for the
accuracy of the model.


% ----------------------------------------------------------------------------
% REPLACES 5.5.7 "Calorie Compensation (Module 8)"
% ----------------------------------------------------------------------------
\subsubsection{Calorie Compensation (Module 8)}
The rule is kept in one function in the app (\texttt{compensation.ts}):
\[
\text{target} = \text{base} + \text{carry-over} + \text{burned today}
\]
The carry-over is minus the total surplus of the last three logged days divided by
three. The surplus of one day is the calories eaten minus (base goal + calories
burned that day). Days with no saved meals are ignored. The carry-over is limited
to $\pm 15\%$ of the base goal, and the final target is never below 1200~kcal or
the user's own base goal if that is lower. Example: with a base goal of 2000~kcal, a
day with 2600~kcal gives a surplus of 600~kcal and lowers the next day's goal by
200~kcal.

Calories burned come from activities that the user logs by hand (walking, running,
cycling, gym, swimming or sports) and a duration. The estimate is
$\text{kcal} = \text{MET} \times \text{weight (kg)} \times \text{hours}$, so 30 minutes
of walking (MET 3.5) for a 70~kg user adds 123~kcal. These values are estimates and
are not measured by a sensor; automatic activity data is planned with the Health
Data Sync module. The home screen shows the adjusted goal, the reason for it and the
disclaimer, and activities are stored in the user's \texttt{activity\_logs} collection
in Firestore.

\textit{Difference from the original design.} The first design divided a weekly
surplus over the remaining days of the week. The built version uses a rolling three-day
window instead, because it reacts the next day and needs no week boundary.


% ----------------------------------------------------------------------------
% ADD to 5.6.3 "History and Dashboard (Modules 2 and 8)"
% ----------------------------------------------------------------------------
The Trends screen also shows a \textbf{Macro Split} card. It adds up protein, carbohydrate
and fat for the chosen range (7, 30 or 90 days) and shows each as a share of the total
calories, using 4, 4 and 9~kcal per gram, with the average grams per logged day. When
no meals are saved in the range the card says so and shows no numbers.


% ----------------------------------------------------------------------------
% ADD a new subsection near 5.5.4 (Recognition Pipeline) for Module 9, and
% relabel the asyncio.gather listing in 5.5.4 as "target design, not built".
% ----------------------------------------------------------------------------
\subsubsection{Live Recognition (Module 9)}
Live recognition is implemented in the app (\texttt{useLiveRecognition.ts}) and uses
the existing recognition endpoint with the YOLOv8 engine. The loop is: capture a
frame at quality 0.3, upload it, show the answer, wait 300~ms and repeat.
In a test over the public server address from a laptop, one request took about
0.77~s, which gives roughly one result per second. A test with a full 1~MB photo
from another computer took about 2.2~s, still inside the 3-second target.
The server does not yet run YOLOv8 and MiDaS at the same time in a thread pool
as described in the design of Section~5.5.4. Only the database calls are
asynchronous, and the depth service is a separate endpoint. Parallel execution is
future work.


% ----------------------------------------------------------------------------
% REPLACES the P2 Mid rows of Table 6.3 "Model Evaluation Targets and Measured
% Values" (keep the other rows).
% ----------------------------------------------------------------------------
\begin{tabular}{|p{5cm}|p{3.2cm}|p{5.2cm}|}
\hline
\textbf{Measure} & \textbf{Target} & \textbf{Measured} \\
\hline
Top-1 accuracy (217 classes, official evaluation) & At least 85\% (Chapter 1) & 58.5\% (top-5: 86.6\%), classification model; target not met \\
\hline
Accuracy change after INT8 quantisation (static) & Small loss only & $-3.7$ points top-1, $-0.2$ points top-5 (comparison set, Table~\ref{tab:quant}) \\
\hline
Model size, original vs.\ static INT8 & Smaller & 6.89~MB vs.\ 1.89~MB (3.6$\times$) \\
\hline
Inference time, original vs.\ static INT8 (CPU, one image) & Within the 3-second goal & 2.4~ms vs.\ 2.1~ms (model only) \\
\hline
\end{tabular}


% ----------------------------------------------------------------------------
% REPLACES the Status column for TC21 - TC24 in Table 6.2
% ----------------------------------------------------------------------------
% TC21 (FR20, FR23)  Passed in part: the adjusted goal and explanation are shown
%                    on an iPhone simulator; the carry-over across two real days
%                    was checked by calculation only.
% TC22 (FR21)        Passed: static INT8 gives top-1 0.663 against 0.700 and the
%                    same speed (Table~\ref{tab:quant}); FP16 not built; the model
%                    is not yet running on the phone.
% TC23 (FR22)        Passed in part on a simulator: the Macro Split card matched
%                    the saved meal (27/34/39\%) and showed its empty state; the
%                    existing calorie, weight and water charts were not part of
%                    this test run.
% TC24 (FR24)        Partly met: live frames are sent one at a time and the app
%                    stays responsive; two simultaneous server requests were not
%                    tested. Live mode on a real camera is not yet tested.


% ----------------------------------------------------------------------------
% ADD to 6.11 "Future Work"
% ----------------------------------------------------------------------------
\begin{itemize}
  \item Run the static INT8 model on the phone for live recognition without the server.
  \item Run YOLOv8 and MiDaS in parallel on the server, as in the original design.
  \item Replace the manual activity entry with automatic data from Google Fit and
        Apple Health (Health Data Sync).
  \item Retrain the classifier to reach the 85\% target, and rebuild the test split
        so the accuracy of the quantised model can be reported on unseen images.
\end{itemize}
